\section{Signed Intelligence Templates and Production Qualification}
\label{sec:aion-flow-templates-production}

This section records the completion of the AION Flow template ecosystem and the locally verifiable
part of production qualification. It is intended to be appended to the existing Pilot and AION
technical document. It does not replace earlier architecture sections.

\subsection{Purpose and trust boundary}

A template is reusable workflow logic, not a copy of a customer's brain. A template package must
therefore exclude credentials, private memory, Business Map facts, account data and granted
authority. Customer state is attached only after installation through local bindings retained by
the customer's mother brain. A template can describe a capability that may eventually send a
message or update a business system, but installing or loading that template never grants the
permission to do so.

The operating sequence is:

\begin{quote}
Signed template $\rightarrow$ dependency and permission review $\rightarrow$ customer-local vault
binding $\rightarrow$ non-executing draft $\rightarrow$ simulation $\rightarrow$ exact approval
$\rightarrow$ governed capability $\rightarrow$ verification and receipt.
\end{quote}

\subsection{Built-in template gallery}

The existing Workflow Canvas exposes a \textbf{Templates} control. It opens an inspectable gallery
rather than a separate automation product. The initial signed catalogue contains:

\begin{itemize}
  \item Local-first research;
  \item Gemini evidence analysis;
  \item Private business analysis;
  \item Multi-model critic and reconciliation;
  \item Governed email drafting; and
  \item Boardroom decision analysis.
\end{itemize}

Each catalogue entry declares an immutable identifier and version, publisher class, department,
description, dependencies, licences, permissions, required bindings and workflow graph. Its
canonical manifest hash is signed with an Ed25519 publisher identity. Inspection verifies that
signature and presents dependency, licence and permission information before use.

The package validator rejects embedded \texttt{vault://} references, redacted secret placeholders,
credential-bearing packages and customer-data-bearing packages. The catalogue API returns public
metadata only. It never returns a raw key, token or customer record.

\subsection{Visual Vault binding and installation}

Templates that require Gemini, private compute, email or critic models declare named binding slots.
The user selects an active Visual Vault binding for each slot. The canvas and installation record
retain only an opaque reference such as \texttt{vault://model/provider-name}; secret material
remains encrypted on the mother brain.

Selecting \textbf{Load logic into canvas} replaces the visible draft with the signed graph and
marks it as a draft. Selecting \textbf{Install for this organisation} pins the exact signed version,
stores only its opaque binding map and records that execution authority remains false. Execution still
requires normal policy, identity, approval and connector checks.

Private organisation libraries require an owner, administrator or workflow-publisher role. Template
versions are immutable. Updates therefore create and pin a new version while preserving the prior
reference. Rollback swaps the active and previous references. Quarantine, revocation and removal
prevent use without destroying historical evidence. Third-party publishing and commercial
settlement remain disabled pending a separate supply-chain and legal review.

\subsection{Production review and signed collaboration}

The Workflow Canvas also exposes a \textbf{Production} centre. Desktop and tablet surfaces support
authoring review; a narrow mobile surface is explicitly read-only. Production checks validate
workflow identity, node identity uniqueness, edge targets, graph size and the caller's declared
review or editing role. The check performs no workflow execution and no external write.

Commits use optimistic concurrency. An editor supplies the revision from which they began. If the
stored revision has advanced, AION rejects the commit with a revision conflict instead of silently
overwriting another person's work. Accepted revisions retain the workflow hash, base revision,
author identity, role, timestamp, mother-brain public key and Ed25519 signature. A version-difference
service identifies added, removed and changed nodes plus the edge-count difference without executing
either version.

For larger graphs, the Production centre can group a selected set of nodes, focus the canvas on a
named group and save selected internal logic as a reusable subflow. Subflows preserve their own
version and definition hash but inherit no credential or execution authority. The visual revision
review shows added, removed and changed node identities and the edge-count change between two signed
revisions.

\subsection{Encrypted portability}

Workflow movement uses a password-protected envelope. AION normalizes the graph to the current
schema, derives a 256-bit key using scrypt with a random salt, and encrypts the package using
AES-256-GCM with a random nonce. The authenticated header identifies the schema, workflow and
expected graph hash without containing workflow content. Import authenticates and decrypts the
envelope, deterministically migrates legacy edge names, checks the resulting graph hash and restores
the workflow with \texttt{authority\_granted=false}. Migration cannot silently restore approvals,
tokens or external-write authority.

\subsection{Zero-content observability}

Customer-owned operational telemetry contains only the workflow hash, node and edge counts, node
families, bounded status totals, latency buckets, timestamp and telemetry hash. It explicitly marks
that it contains no workflow content, credentials or customer identifiers. This permits local
health and performance review without turning Tessaris into a collector of business prompts,
documents or results.

\subsection{Verified scope and remaining qualification}

Automated tests currently prove signed template verification, state-free packages, opaque binding
enforcement, role-based private publication, version pinning, rollback, quarantine, revocation,
least-privilege review, revision conflict detection, encrypted round trips, zero-content telemetry
and adversarial rejection of malformed or excessive graphs.

Production qualification is not yet complete. The remaining work is:

\begin{enumerate}
  \item complete a formal threat model and third party dependency review;
  \item benchmark parallel routes and long running recovery under representative loads; and
  \item perform accessibility, privacy, regulatory and representative medium business field
        qualification.
\end{enumerate}

These items remain openly recorded because local automated checks cannot substitute for independent
field evidence. Until they are complete, AION Flow may be locally qualified for controlled use but
must not be represented as fully enterprise-qualified.
